\sect{一、单项选择题（每小题 2 分，共 20 分）}
\begin{enumerate}[label=\arabic*.]
\item \judge 某团队已留出一份独立测试集，准备用其余数据做五折交叉验证，选择分类器的超参数。模型使用从数据估计的标准化参数。哪种流程正确？
\optss{先在五折数据的全集拟合标准化参数，再划分每一折进行训练和验证}{每折在验证子集上拟合标准化参数，因为该折的分数要在这个子集上比较}{每折只用训练子集拟合标准化参数，再用它变换该折的训练、验证子集。选定超参数后，用全部非测试数据重新拟合，最后在测试集上评估}{每个候选超参数均在留出测试集上测一次，并用测试分数最高的结果作为最终评估}

\item \judge 某疾病筛查系统把患病记为正类。共检测 $200$ 人，其中 $\mathrm{TP}=72$、$\mathrm{FP}=8$、$\mathrm{FN}=18$、$\mathrm{TN}=102$。准确率、精确率、召回率、$F_1$ 依次为多少？
\opts{$0.8700,\ 0.9000,\ 0.8000,\ 0.8471$}{$0.8700,\ 0.8000,\ 0.9000,\ 0.8471$}{$0.8700,\ 0.9000,\ 0.8000,\ 0.8500$}{$0.8000,\ 0.9000,\ 0.8000,\ 0.8471$}

\item \judge 某二分类模型给两个正样本的分数为 $0.90,0.50$，给两个负样本的分数为 $0.80,0.30$。分数不低于阈值时判正。阈值为 $0.70$ 时的 $(\mathrm{FPR},\mathrm{TPR})$ 和遍历所有阈值得到的 AUC 分别是多少？
\opts{$(0.5,0.5)$，$0.50$}{$(0,0.5)$，$0.75$}{$(0.5,1)$，$0.75$}{$(0.5,0.5)$，$0.75$}

\item \judge 设 $P=(1/2,1/4,1/4)$，$Q=(1/4,1/8,5/8)$。以 $2$ 为对数底，记 $H(P,Q)=-\sum_i P_i\log_2Q_i$。计算 $H(P),H(Q),H(P,Q),\KL(P\,\|\,Q)$，结果保留两位小数。
\opts{$1.50,\ 1.30,\ 1.75,\ 0.25$}{$1.50,\ 1.30,\ 1.92,\ 0.42$}{$1.50,\ 1.30,\ 1.92,\ 0.62$}{$1.30,\ 1.50,\ 1.75,\ 0.45$}

\item \judge 决策树某节点有 $5$ 个正样本、$5$ 个负样本。用某属性分成两支后，两支的正负计数依次为 $(1,3)$、$(4,2)$。以 $2$ 为对数底，增益比取信息增益除以分支比例的熵，基尼指数取两支基尼指数的加权平均。两者依次约为多少？
\opts{$0.1282,\ 0.4167$}{$0.1245,\ 0.4167$}{$0.1282,\ 0.08333$}{$0.1282,\ 0.4097$}

\item \judge 对两个类别做最小错误率分类。各类的类条件密度均为二维高斯，协方差矩阵正定，先验概率非零。下列哪项正确？
\optss{对每个输入只比较两个类条件密度，任意类别先验都适用}{两类协方差相等时，如果类别先验不相等，对数判别函数之差就会含二次项}{比较各类的先验概率与类条件密度的乘积。协方差相等时对数判别函数之差是仿射函数，协方差不等时含二次项}{协方差不等时，取对数会消掉逆协方差矩阵之间的差异}

\item \judge 用 1-NN 给原点分类，训练点为 $(3,0)$、$(2,2)$、$(4,3)$，标签依次为甲、乙、丙。不做特征缩放。采用曼哈顿距离和欧氏距离时，分别判为哪一类？
\opts{乙、乙}{甲、甲}{乙、甲}{甲、乙}

\item \judge 两层网络的某一层有 $64$ 个输入和 $32$ 个输出。分别采用针对 $\tanh$ 的标准 Xavier 初始化（增益取 $1$）和针对 ReLU 的 He 初始化（按扇入设置）。两种零均值初始化的权重方差依次是多少？
\opts{$1/32,\ 1/48$}{$1/48,\ 1/32$}{$1/96,\ 1/64$}{$1/48,\ 1/16$}

\item \judge 实验课准备比较 LeNet 和 AlexNet。下列哪组描述符合这两个经典网络的设计？
\optss{LeNet 用卷积和下采样处理手写字符。AlexNet 在更深的卷积网络中采用 ReLU，并在部分全连接层使用 Dropout}{LeNet 用跨层恒等残差连接训练卷积层。AlexNet 首创多分支 Inception 模块}{LeNet 把图像切成块后使用多头自注意力。AlexNet 依次使用多个门控循环单元}{LeNet 依靠全连接层直接分类，省去卷积。AlexNet 主要用循环状态整合空间信息}

\item \judge 在带偏置的二分类感知机中，令 $\tilde{\bm x}_i=(\bm x_i^\top,1)^\top$、$y_i\in\{-1,+1\}$。增广权重 $\bm w$ 从零开始，学习率为 $1$，遇到 $y_i\bm w^\top\tilde{\bm x}_i\le0$ 时更新。设 $m_i$ 为样本 $i$ 触发更新的次数。若 $\|\tilde{\bm x}_i\|\le R$ 且有单位向量 $\bm u$ 满足 $y_i\bm u^\top\tilde{\bm x}_i\ge\gamma>0$，哪项正确？
\optss{对偶展开为 $\bm w=\sum_i m_i\tilde{\bm x}_i$，总更新次数至多 $(R/\gamma)^2$}{对偶展开为 $\bm w=\sum_i m_i y_i\tilde{\bm x}_i$，对线性不可分的训练集也保证有限次停止更新}{对偶展开为 $\bm w=\sum_i m_i y_i\tilde{\bm x}_i$，总更新次数至多 $(R/\gamma)^2$}{每次只增加一个计数，所以总更新次数至多为训练样本数}
\end{enumerate}

\sect{二、判断对错题（每小题 1 分，共 10 分）}
\begin{enumerate}[label=\arabic*.,start=11]
\item \judge 在固定输入 $x$ 处，以平方误差衡量回归预测，训练集随机抽取，测试对象独立于训练集。期望预测误差等于估计量相对 $\E(Y\mid x)$ 的偏差平方、估计量自身的方差和 $\operatorname{Var}(Y\mid x)$ 三项之和。
\item \judge 一组带标签的二维数据沿某个低方差方向能很好地区分两类，所以只保留一个方向的 PCA 一定会选择该方向，与 LDA 的选择相同。
\item \judge 用验证集选择训练轮次属于早停。若裁剪后的图片仍能辨认原物体且类别不变，可以把裁剪结果用于训练数据增强。单独执行的权重衰减可以在梯度更新之外，将权重乘以小于 $1$ 的正因子。
\item \judge 对 $f(w)=w^2/2$ 从非零 $w$ 开始做梯度下降，学习率分别取 $0.01$ 和 $2.2$。两种设置下 $|w|$ 都逐步减小，取 $2.2$ 时减小得更快。
\item \judge 当链式反传每一步的标量导数绝对值均为 $0.5$ 或均为 $1.5$ 时，其连乘分别会随层数增加而趋零或增长。按范数裁剪可限制过大的梯度，但不会放大因连乘而衰减的小梯度。
\item \judge 对常见的多层 RNN，同一时刻上一循环层的输出可作为下一循环层的输入。若只看输入输出序列长度，给定一张图像生成描述、读完整句子判情感、逐词标注词性分别可建为一对多、多对一、等长多对多。
\item \judge 在标准 VAE 中，若解码器对给定样本的重构对数似然与 $z$ 无关，那么对该样本最大化 ELBO 会促使近似后验远离先验，从而保证隐变量保留样本信息。
\item \judge 用 one-hot 标签计算 softmax 交叉熵时，损失式只显式出现正确类的对数概率，因此提高任何一个错误类的 logit 都不改变该样本的损失。
\item \judge 对所有样本的实值预测分数施加同一个严格递增变换，并对变换后的分数扫描全部阈值，所得 ROC 曲线和 AUC 都保持不变。
\item \judge 在高斯朴素贝叶斯中，若某类下指定特征值的条件密度大于 $1$，说明所用的高斯分布无效。
\end{enumerate}

\sect{三、计算与分析题（共 8 题，共 70 分）}

\noindent 如无另行规定，$\log$ 表示自然对数，计算结果可保留精确式或写到 4 位有效数字。

\noindent\textbf{1. 多层感知机与反向传播}\pts{10}

某设备根据三个传感读数预测故障概率。网络只有一层隐藏层，记 $r=\tfrac12\ln 3$，输入及参数为
\[
\bm{x}=(1,-1,2)^{\mathsf T},\quad
\bm{W}=\begin{pmatrix}r&0&0\\0&r&0\end{pmatrix},\quad \bm{b}=\bm{0},
\qquad \bm{v}=(2,-1)^{\mathsf T},\quad c=\ln 3-\tfrac32.
\]
隐藏层计算 $\bm{z}=\bm{W}\bm{x}+\bm{b}$、$\bm{h}=\tanh(\bm{z})$，其中 $\tanh$ 逐元素作用。输出层计算 $a=\bm{v}^{\mathsf T}\bm{h}+c$、$p=\sigma(a)$。标签 $y=1$，损失为二元交叉熵 $L=-y\ln p-(1-y)\ln(1-p)$，对数以 $e$ 为底。各梯度在更新前计算。

\begin{enumerate}[label=(\arabic*)]
\item 计算隐藏层的 $\bm{z}$、$\bm{h}$。\pts{1}
\item 计算输出得分 $a$、概率 $p$ 和损失 $L$。\pts{1}
\item 从二元交叉熵与 Sigmoid 的导数出发，推导 $\partial L/\partial a$，并计算其值。\pts{1}
\item 计算 $L$ 对输出层参数 $\bm{v},c$ 的梯度。\pts{1}
\item 用链式法则计算隐藏层误差信号 $\partial L/\partial\bm{z}$。\pts{1}
\item 计算 $L$ 对隐藏层参数 $\bm{W},\bm{b}$ 的梯度。\pts{1}
\item 对上述梯度做第一次 Adam 更新。取一阶、二阶矩的衰减系数 $\beta_1=1/2$、$\beta_2=3/4$，初始矩均为零。以 $c$ 的梯度为例，求更新后的一阶矩、二阶矩及各自偏差修正后的值。\pts{1}
\item Adam 的学习率为 $1/20$，稳定项 $\epsilon=1/16$ 加在二阶矩估计的平方根外。求这一步更新后的 $\bm{W},\bm{b}$。\pts{1}
\item 求同一步 Adam 更新后的 $\bm{v},c$。\pts{1}
\item 若改为对所有原始参数梯度一起做全局 $\ell_2$ 范数裁剪，阈值为 $1$，求统一缩放因子和裁剪后的 $c$ 梯度。\pts{1}
\end{enumerate}

\noindent\textbf{2. 支持向量机与主成分特征}\pts{8}

四个二维训练样本依次为 $\bm x_1=(-2,-4)^\top$、$\bm x_2=(-2,0)^\top$、$\bm x_3=(0,2)^\top$、$\bm x_4=(4,2)^\top$，对应标签 $(-1,-1,+1,+1)$。先对样本做 PCA，保留第一主成分，再用投影坐标训练一维 SVM。协方差矩阵按 $1/4$ 计算。第一主轴取两个分量均为正的单位向量，投影坐标记作 $z_i$。

\begin{enumerate}[label=(\arabic*)]
\item 求样本均值和协方差矩阵 $S$。\pts{1}
\item 求 $S$ 的两个特征值、第一主轴和四个投影坐标。\pts{1}
\item 对任意单位向量 $\bm u$，写出投影方差，用拉格朗日乘子法推出使它取极值的特征值方程。说明为何应选最大特征值对应的方向。\pts{1}
\item 在投影坐标上，写出硬间隔约束，求最优判别函数 $f(z)=wz+b$ 与两条间隔边界之间的宽度。\pts{1}
\item 对硬间隔问题写出关于 $w,b$ 的两个拉格朗日平稳条件，乘子记作 $\alpha_i$。\pts{0.5}
\item 求四个硬间隔对偶乘子，并根据间隔约束和互补松弛说明哪些点是支持向量。\pts{1}
\item 对新点 $(-1,-2)^\top$ 计算 $f$ 值并预测标签。\pts{0.5}
\item 用同一训练集改做软间隔 SVM，取 $C=1/8$。给出候选对偶乘子 $(0,1/8,1/8,0)$。检查盒约束和等式约束，并计算对应的 $w$。\pts{1}
\item 对（8）取 $b=0$，求最小松弛变量，指出哪些点进入间隔内部。\pts{0.5}
\item 考虑作用于实数投影坐标的函数 $K(z,z')=(zz')^2$。写出一个显式特征映射 $\varphi$，使 $K(z,z')=\varphi(z)\varphi(z')$。\pts{0.5}
\end{enumerate}

\noindent\textbf{3. 二维测点的密度估计与聚类}\pts{12}

某装置记录了五个二维测点。分类时用表中的标签，聚类时只用坐标。记查询点为 $q=(0,0)$。
\[
\begin{array}{c|ccccc}
\text{测点}&a&b&c&d&e\\\hline
\text{坐标}&(1,0)&(1,1)&(0,3)&(2,2)&(1,4)\\
\text{标签}&+&-&-&+&-
\end{array}
\]
\begin{enumerate}[label=(\arabic*)]
\item 用各向同性高斯核估计 $q$ 的密度。核为 $K_h(u)=(2\pi h^2)^{-1}\exp[-\|u\|_2^2/(2h^2)]$，带宽满足 $h^2=1/(2\ln 2)$。写出密度估计式并代入数据。保留精确表达式即可。\pts{1}
\item 取 $k=3$，分别按曼哈顿距离和欧氏距离列出 $q$ 的前三个近邻，并用不加权投票给出预测标签。\pts{1}
\item 仪器又记录了三个与 $e$ 坐标完全相同的测点。把它们都视为独立观测，保持带宽不变。用第 (1) 问的旧估计 $\widehat p(q)$ 和核函数 $K_h$ 表示更新后的 Parzen 密度，并判断它比原估计大还是小。\pts{1}
\item K-means++ 已选 $a,d$ 为两个中心。第三个中心落在 $\{b,e\}$ 中的概率是多少？已知它落在 $\{b,e\}$ 中，求它就是 $e$ 的条件概率。\pts{1}
\item 以 $a,d$ 为两个初始中心，暂不加入第 (4) 问讨论的第三个中心。用平方欧氏距离对表中的五个样本进行一次 K-means 分配，列出两个簇和更新前的簇内平方误差和。\pts{1}
\item 用第 (5) 问的分配同时更新两个中心，求新中心以及按该分配计算的簇内平方误差和。\pts{1}
\item 另有四个对象 $U,V,W,Z$，距离如下。用单链法做凝聚层次聚类，画出完整树状图并标出每次合并的高度。
\[
\begin{array}{c|rrrr}
 &U&V&W&Z\\\hline
U&0&4&10&19\\
V&4&0&6&15\\
W&10&6&0&9\\
Z&19&15&9&0
\end{array}
\]\pts{1}
\item 对第 (7) 问的距离矩阵，分别画出全链法和平均链法的完整树状图，标出合并高度。平均链按两簇之间所有跨簇点对的距离取平均。\pts{2}
\item 对表中的五个二维测点运行 DBSCAN。用欧氏距离，$\varepsilon=\sqrt{2}$，$\mathrm{MinPts}=3$。邻域包括自身及距离等于 $\varepsilon$ 的点。列出核心点、边界点和噪声点。\pts{2}
\item 若用带固定正权重、共同球形协方差 $\sigma^2 I$ 的 GMM 代替 K-means，均值更新时如何使用 E 步的责任度？当 $\sigma^2$ 趋于零且每个样本的最近中心唯一时，解释软分配如何趋向硬分配。\pts{1}
\end{enumerate}

\noindent\textbf{4. 自编码器、证据下界与高斯变分后验}\pts{8}

先考虑一个离散生成模型。潜变量 $z$ 取 $0$ 或 $1$，先验各为 $1/2$。已知 $p(x=1\mid z=0)=3/4$、$p(x=1\mid z=1)=1/4$。观测到 $x=1$，变分后验取 $q(0\mid x)=q(1\mid x)=1/2$。各对数均以 $e$ 为底。

\begin{enumerate}[label=(\arabic*)]
\item 对同一幅图像，分别说明普通确定性自编码器和 VAE 的编码器输出什么。用 VAE 生成新图像时，潜变量从哪个分布抽取？\pts{0.5}
\item 求 $p(x=1)$，并求两个潜在状态对应的比值 $p(x=1,z)/q(z\mid x=1)$。\pts{1}
\item 当 $q$ 在两个状态都为正时，把 $p(x)$ 写成关于 $q$ 的期望，再用 Jensen 不等式推导对数证据的下界。\pts{1.5}
\item 用 $p(x,z)=p(z)p(x\mid z)$ 将上问下界拆成重构项和先验 KL 项。\pts{1}
\item 求给定 $q$ 下的 ELBO，以及 $\ln p(x=1)$ 与 ELBO 的差值。\pts{1}
\item 用 Jensen 等号条件求出使下界等于对数证据的 $q(0\mid x=1)$ 和 $q(1\mid x=1)$。\pts{1}
\item 改用一维连续模型，设 $q_c(z)=\mathcal N(\mu,v)$、$p_c(z)=\mathcal N(0,1)$，其中 $v>0$ 是方差。写出两者的对数密度比 $\ln[q_c(z)/p_c(z)]$。\pts{0.5}
\item 用 $q_c$ 下的二阶矩对上述比值的对数取期望，推导 $\KL(q_c\Vert p_c)$ 的闭式结果。\pts{1}
\item 回到离散模型。若改用只在 $z=1$ 取值的 $q$，则 $\E_q[p(x=1,z)/q(z\mid x=1)]$ 还等于 $p(x=1)$ 吗？说明原因。\pts{0.5}
\end{enumerate}

\noindent\textbf{5. 卷积、逆向序列网络与掩码注意力}\pts{8}

\begin{enumerate}[label=(\arabic*)]
\item 一张特征图的尺寸为 $5\times6\times2$，依次为高、宽、通道数。用 $3$ 个空间尺寸为 $2\times2$ 的卷积核做互相关，两个空间方向的步长均为 $2$，上下左右各填充 $1$ 个像素宽的零，每个输出通道有一个偏置。写出输出高、宽的计算式及输出尺寸。\pts{1}
\item 第一个输出通道在某位置看到的两个输入通道块依次为 $\begin{pmatrix}1&0\\2&1\end{pmatrix}$、$\begin{pmatrix}0&2\\1&-1\end{pmatrix}$，对应核切片为 $\begin{pmatrix}1&-1\\0&1\end{pmatrix}$、$\begin{pmatrix}0&1\\-1&0\end{pmatrix}$，偏置为 $1/2$。求该位置的输出，并求整个第（1）问卷积层的可训练参数量。\pts{1}
\item 在第（1）问输出之后接一个步长为 $1$、不填充、有偏置的 $1\times1$ 卷积层，输出 $4$ 个通道。求输出尺寸和参数量。\pts{0.5}
\item 设同形残差块为 $\bm{y}=\bm{x}+F(\bm{x})$。写出输入、输出之间的雅可比矩阵，说明梯度从哪条路径直接传回，以及新增块如何保留浅层网络已学到的映射。\pts{0.5}
\item 逆向 RNN 按 $t=3,2,1$ 的顺序计算，$\bm{h}_t=\bm{W}\bm{h}_{t+1}+\bm{x}_t$，$\bm{h}_4=\bm{0}$。当前 $\bm{W}=\operatorname{diag}(1/2,1/3)$，$\bm{x}_1=(0,1)^{\mathsf T}$、$\bm{x}_2=(1,0)^{\mathsf T}$、$\bm{x}_3=(2,-1)^{\mathsf T}$。输出层对权重 $\bm{v}=(1,0)^{\mathsf T}$ 的线性得分做 Sigmoid。写出 $o_t$ 的公式，求 $\bm{h}_2,o_2$。\pts{1}
\item 三个时刻的损失相加得到 $L$。只计本时刻损失时，三个损失对各自隐状态的梯度依次是 $(1,0)^{\mathsf T}$、$(0,1)^{\mathsf T}$、$(0,0)^{\mathsf T}$。写出 BPTT 汇总梯度的递推方向和公式，再求 $L$ 对共享权重 $W_{11}$ 的偏导。\pts{1}
\item 对 LSTM，使用标准细胞状态更新 $\bm{c}_t=\bm{f}_t\odot\bm{c}_{t-1}+\bm{i}_t\odot\widetilde{\bm{c}}_t$。分析某一维的细胞状态直接传递路径时，固定门值和候选值。写出沿该路径跨 $k$ 步的梯度因子，并说明什么样的遗忘门值有利于长期传递。\pts{1}
\item 某多头自注意力的输入是长度为 $n$、每位置维数为 $8$ 的序列。两个头的查询、键、值各为 $4$ 维，拼接后再投影到 $8$ 维。各投影均无偏置。给出每头的三个投影矩阵、拼接与输出投影的尺寸，以及总参数量。\pts{1}
\item 某单头的缩放点积得分矩阵是下列矩阵逐元素取自然对数后的结果
$\begin{pmatrix}2&3&1\\2&1&3\\1&3&2\end{pmatrix}$。
行是查询位置，列是键位置，允许每个位置看自身及更早的位置。写出应用加性因果掩码和逐行 Softmax 后的完整权重矩阵。\pts{0.5}
\item 若各头维数总和为 $d$，说明为何稠密实现中，仅计算注意力得分和加权求和就需要 $O(n^2d)$ 时间。\pts{0.5}
\end{enumerate}

\noindent\textbf{6. 工单分类中的朴素贝叶斯}\pts{8}

某设备厂将工单分为故障类 $F$ 和正常类 $N$，两类先验均为 $1/2$。三种特征记录方式的数据如下。每种方式单独建模，每类各有四条训练工单。类别特征统计如下。
\[
\begin{array}{c|cc}
 &\text{紧急程度（急、普、例行）}&\text{地点（室内、室外）}\\\hline
F&(3,1,0)&(1,3)\\
N&(1,2,1)&(3,1)
\end{array}
\]
词表按顺序为“振动、发热、异响”，各词在训练文本中的统计如下。
\[
\begin{array}{c|cc}
 &\text{累计词频}&\text{含该词的文档数}\\\hline
F&(5,2,1)&(3,2,1)\\
N&(1,2,5)&(1,2,3)
\end{array}
\]
各类在两个连续传感器上的读数如下。
\[
\begin{array}{c|rrrr}
F&(2,1)&(2,3)&(4,1)&(4,3)\\
N&(0,-1)&(0,3)&(2,-1)&(2,3)
\end{array}
\]
\begin{enumerate}[label=(\arabic*)]
\item 将紧急程度和地点作为类别特征。对每个特征的各水平分别作加一平滑，求两类下“急”和“室外”的条件概率。\pts{1}
\item 新工单为“急、室外”。用第 (1) 问建立的类别特征朴素贝叶斯模型求后验概率 $P(F\mid\text{急、室外})$，并给出分类。\pts{1}
\item 对工单文本采用多项式朴素贝叶斯，词频作加一平滑。写出两类的词条件概率向量，并说明分母如何确定。\pts{1}
\item 一条新工单的文本为“振动、振动、发热”。用第 (3) 问模型计算 $F$ 相对 $N$ 的后验对数优势 $\ln(P(F\mid x)/P(N\mid x))$，并求 $P(F\mid x)$。\pts{1}
\item 仍对第 (4) 问的文本，改用加一平滑的伯努利朴素贝叶斯。求两类下该文档的条件概率与分类，明确列出未出现的词贡献的因子。\pts{1}
\item 改用连续读数。对两类的两个特征各自用极大似然法估计高斯均值和方差，给出两组均值向量及方差向量。\pts{1}
\item 新工单的连续读数为 $(2,2)$。用第 (6) 问的高斯朴素贝叶斯模型求它的后验对数优势、故障后验概率和预测类别。对数用自然对数。\pts{1}
\item 说明朴素贝叶斯和逻辑回归分别对哪些概率建模。若文本有很多词，解释用对数计算未归一化类别得分的作用。\pts{1}
\end{enumerate}

\noindent\textbf{7. 传感器校准、回归与正则化}\pts{8}

用三个校准样本 $(-1,1),(0,1),(1,3)$ 训练回归模型，第一列是输入 $x$，第二列是目标 $y$。验证集为 $(-2,1),(2,4)$，只用来比较模型。线性模型为 $w_0+w_1x$，二次模型为 $v_0+v_1x+v_2x^2$，均用最小二乘拟合。MSE 与 MAE 均除以计算所用的样本数。

\begin{enumerate}[label=(\arabic*)]
\item 写出线性模型训练时的设计矩阵 $X$ 与目标向量 $\bm y$。\pts{0.5}
\item 写出正规方程，算出左侧的 $X^\top X$ 和右侧的 $X^\top\bm y$。\pts{1}
\item 求线性模型的系数并写出预测函数。\pts{1}
\item 求二次模型的三个系数。\pts{1}
\item 分别求两个模型在验证集上的 MSE 与 MAE。\pts{1}
\item 哪个模型在给定验证集上表现更好？结合训练误差和验证误差解释为何不能只按训练误差选阶数。\pts{0.5}
\item 假设各校准样本的噪声独立，服从零均值、固定方差 $\sigma^2>0$ 的高斯分布。写出线性模型的负对数似然中依赖 $\bm w$ 的部分，并解释它与最小二乘的关系。\pts{1}
\item 另用逻辑回归预测传感器是否报警。某样本的 logit 为 $\ln 2$。求报警概率，并给出以 $0.75$ 为阈值时的分类结果（概率大于等于阈值判报警）。\pts{0.5}
\item 对无正则化的二元逻辑回归，若训练集严格线性可分，设某组参数给所有样本的有符号间隔 $m_i=y_i(\bm w^\top\bm x_i+b)>0$，其中 $y_i\in\{-1,1\}$。将整组参数乘 $t>0$，写出总交叉熵损失对 $t$ 的表达式，说明参数范数与损失如何变化。\pts{1}
\item 对（9）中的所有参数加入平方范数惩罚，系数固定为正。说明为什么继续沿这条射线放大参数不再能使目标任意降低。\pts{0.5}
\end{enumerate}

\noindent\textbf{8. 滤器状态的识别与更新}\pts{8}

用 HMM $\lambda=(A,B,\pi)$ 描述滤器。隐藏状态 1 表示洁净，状态 2 表示积尘。读数分为低、中、高三个等级，记为 $L,M,H$。$A$ 的行、列依次是当前状态和下一状态，$B$ 的列按 $L,M,H$ 排列。参数及观测序列为
\[
 A=\frac14\begin{pmatrix}3&1\\1&3\end{pmatrix},\qquad
 B=\frac14\begin{pmatrix}2&1&1\\1&1&2\end{pmatrix},\qquad
 \pi=(1/2,1/2),\qquad O=(L,M,H,L).
\]
记时刻 $t$ 的状态为 $S_t$，观测为 $o_t$。以下均可保留分数。
\begin{enumerate}[label=(\arabic*)]
\item 写出前向量 $\alpha_t(i)=P(o_1,\ldots,o_t,S_t=i\mid\lambda)$ 的初始化式和递推式，指出转移概率项与发射概率项。\pts{1}
\item 计算 $\alpha_1$。\pts{0.5}
\item 计算 $\alpha_2$。\pts{0.5}
\item 计算 $\alpha_3$。\pts{0.5}
\item 计算 $\alpha_4$。\pts{0.5}
\item 计算 $P(O\mid\lambda)$。\pts{0.5}
\item 记 $\beta_t(i)=P(o_{t+1},\ldots,o_4\mid S_t=i,\lambda)$。写出 $\beta_4$ 与后向递推式，再计算 $\beta_3$。\pts{0.5}
\item 继续计算 $\beta_2$ 和 $\beta_1$。\pts{0.5}
\item 求时刻 2 的两个状态后验 $\gamma_2(i)=P(S_2=i\mid O,\lambda)$。\pts{0.5}
\item 记 $\xi_t(i,j)=P(S_t=i,S_{t+1}=j\mid O,\lambda)$。写出 $\xi_t(i,j)$ 的前后向计算式，并算出 $\xi_t(1,2)$ 在 $t=1,2,3$ 时的值。\pts{1}
\item 接着做一轮 Baum--Welch 更新。写出初始概率、转移概率、发射概率的一般更新式，并算出更新后的 $a_{12}^{\mathrm{new}}$。\pts{2}
\end{enumerate}

\vfill \begin{center}\zihao{5} —— 试题结束 ——\end{center}
